\documentclass{article}
\usepackage{geometry}
\geometry{a4paper, margin=1in}
\title{GCLI.md}
\author{}
\date{}

\begin{document}
\maketitle
\section*{all\_demos.py}
\subsection*{run\_all\_demos()}
Runs all demo functions, and either shows them inline or saves them to a single HTML file.

\section*{common.py}
\subsection*{initialize\_standard\_simulation(start\_time: datetime)}
Initializes a standard simulation with a predefined set of satellites. It consolidates TLE data, initializes the main data structure, and propagates satellites to their initial positions.

\section*{constants.py}
This file defines global constants used as indices for NumPy arrays throughout the simulation. This avoids using "magic numbers" and makes the code more readable. It includes:
\begin{itemize}
\item Radii for Earth and Moon in meters.
\item Detector array indices for properties like aperture, pixel size, quantum efficiency, etc.
\item Orbital elements array indices for semi-major axis, eccentricity, inclination, etc.
\item Pointing state array indices for pointing count and place.
\end{itemize}

\section*{demo1.py}
\subsection*{demo1()}
Runs a full demonstration of the simulation tools. It generates and returns a Plotly figure object for the satellite positions plot.

\section*{demo2.py}
\subsection*{demo2()}
Runs a demonstration plotting satellite and celestial positions at two different times. It returns a Plotly figure object.

\section*{demo3.py}
\subsection*{demo3()}
Runs a demonstration plotting a single LEO satellite trajectory over 90 minutes. It returns a Plotly figure object.

\section*{demo4.py}
\subsection*{demo4()}
Runs a demonstration plotting a single GEO satellite trajectory over 23 hours. It returns a Plotly figure object.

\section*{demo\_exclusion\_debug\_print.py}
\subsection*{demo\_exclusion\_debug\_print()}
Demonstrates the debug printing feature of the exclusion function. It runs an exclusion check and prints detailed debug output.

\section*{demo\_exclusion\_table.py}
\subsection*{demo\_exclusion\_table()}
Demonstrates the creation and visualization of the exclusion table. It generates and returns a Plotly heatmap figure.

\section*{demo\_fixedpoints.py}
\subsection*{demo\_fixedpoints()}
Demonstrates the fixedpoints data structure by plotting it in 3D. It returns a Plotly figure object.

\section*{demo\_lambertian.py}
\subsection*{demo\_lambertian()}
Runs a demonstration of the lambertiansphere function, including example calculations and a plot of brightness vs. phase angle.

\section*{demo\_pointing\_plot.py}
\subsection*{demo\_pointing\_plot()}
Demonstrates the plot\_pointing\_vectors function. It returns a Plotly figure object showing satellite positions with pointing vectors.

\section*{demo\_pointing\_sequence.py}
\subsection*{demo\_pointing\_sequence()}
Demonstrates the satellite pointing sequence functionality. It animates the pointing vectors of satellites over several time steps.

\section*{demo\_pointing\_vectors.py}
\subsection*{demo\_pointing\_vectors()}
Demonstrates the generation and plotting of pointing vectors. It generates a number of vectors and plots them on a sphere.

\section*{demo\_sky\_scan.py}
\subsection*{demo\_sky\_scan()}
Performs a sky scan from a GEO satellite to map celestial exclusion zones. It generates a heatmap of the sky showing clear and excluded areas.

\section*{show\_geo\_search.py}
\subsection*{show\_geo\_search()}
A demo that performs a geometric search for satellites using a GEO constellation, and generates several plots to visualize the satellite pointing updates and the RA/Dec history of one satellite.

\section*{generate\_log\_spherical\_points.py}
\subsection*{generate\_log\_spherical\_points(num\_points: int, inner\_radius: float, outer\_radius: float, object\_size\_m: float = 1.0, seed: int = None)}
Generates 3D points with logarithmic radial and uniform angular distribution. It creates a point cloud for use as fixed points in the simulation.

\section*{generate\_report.py}
\subsection*{generate\_demo\_html\_report()}
Runs all plotting demos and saves the output to a single HTML file named demo\_plots.html.

\section*{lambertian.py}
\subsection*{lambertianSphere(vec\_from\_sphere\_to\_light: np.ndarray, vec\_from\_sphere\_to\_observer: np.ndarray, albedo: float, radius: float)}
Calculates the effective brightness of a Lambertian sphere. It determines the apparent brightness based on the angle between the light source and the observer.

\section*{plotting\_3d.py}
\subsection*{plot\_3d\_scatter(positions: np.ndarray, title: str, plot\_time: datetime, labels: Optional[List[str]] = None, marker\_size: int = 1, trace\_name: str = 'Points')}
Creates a 3D plot of object positions. It can plot any set of 3D points and includes a representation of the Earth.

\section*{plotting\_vectors.py}
\subsection*{plot\_pointing\_vectors(data\_struct: Dict[str, Any], title: str, plot\_time: datetime)}
Creates a 3D plot of satellites with their pointing vectors. It shows the direction each satellite is pointing.

\section*{pointing.py}
\subsection*{pointing\_place\_update(data\_struct: Dict[str, Any])}
Increments the pointing place for all satellites, wrapping around if necessary.
\subsection*{jerk(data\_struct: Dict[str, Any], satellite\_number: int)}
Moves the pointing vector of a specific satellite by 0.3 radians in a random direction.
\subsection*{generate\_pointing\_sphere(data\_struct: Dict[str, Any], n\_points: int)}
Generates a pointing sphere with n\_points and stores it in the data\_struct.
\subsection*{update\_satellite\_pointing(data\_struct: Dict[str, Any])}
Updates the pointing vector for each satellite based on its pointing state.
\subsection*{find\_and\_jerk\_blind\_satellites(data\_struct: Dict[str, Any])}
Finds satellites with no visibility and applies the 'jerk' function to them.

\section*{pointing\_vectors.py}
\subsection*{pointing\_vectors(n: int)}
Generates n equally spaced points on a unit sphere using the Fibonacci lattice algorithm.
\subsection*{plot\_vectors\_on\_sphere(vectors: np.ndarray, title: str)}
Creates a 3D plot of vectors on a sphere.

\section*{propagation.py}
\subsection*{celestial\_update(data\_struct: Dict[str, Any], time\_date: datetime)}
Calculates and updates the positions of the Sun and Moon for a given time.
\subsection*{readtle(tle\_file\_path: str)}
Reads a TLE file and extracts orbital elements and epochs for each satellite.
\subsection*{propagate\_satellites(data\_struct: Dict[str, Any], time\_date: datetime)}
Updates satellite positions and pointing vectors based on their orbital elements to a specified time.

\section*{radiometry\_calcs.py}
\subsection*{mag(x: float)}
Calculates a magnitude value from a linear ratio.
\subsection*{amag(x: float)}
Calculates the linear ratio from a magnitude value.
\subsection*{\_planck\_law(wav\_m: float, temp\_k: float)}
Helper function for Planck's law for spectral radiance.
\subsection*{blackbody\_flux(temperature: float, lambda\_short: float, lambda\_long: float)}
Numerically computes the integrated spectral radiance of a blackbody over a given wavelength band.
\subsection*{stefan\_boltzmann\_law(temperature: float)}
Calculates the total power radiated per unit area by a blackbody.
\subsection*{plot\_blackbody\_spectrum(temperature: float)}
Plots the spectral radiance of a blackbody from 0.5 to 30 microns.
\subsection*{plot\_blackbody\_spectrum\_visible\_nir(temperature: float)}
Plots the spectral radiance of a blackbody from 0.1 to 1 micron.

\section*{radiometry\_data.py}
This file contains radiometric data and physical constants.
\begin{itemize}
\item \texttt{AU\_M}: Astronomical Unit in meters.
\item \texttt{RSUN\_M}: Radius of the Sun in meters.
\item \texttt{FILTER\_DATA}: A dictionary containing data for standard astronomical filters (Johnson-Cousins, NIR, SDSS, Mid-IR, JWST MIRI).
\begin{itemize}
\item For each filter, it provides:
\begin{itemize}
\item \texttt{sun}: Apparent magnitude of the sun.
\item \texttt{sky}: Sky brightness in magnitudes per square arcsecond.
\item \texttt{central\_wavelength}: Central wavelength in nanometers.
\item \texttt{bandwidth}: Bandwidth in nanometers.
\item \texttt{zero\_point}: Photon flux for a 0-magnitude object.
\end{itemize}
\end{itemize}
\end{itemize}

\section*{simulation.py}
\subsection*{initializeStructures(num\_satellites: int, num\_observatories: int, num\_red\_satellites: int, start\_time: datetime, delta\_time: float = 60.0)}
Initializes a minimal, empty data structure for a space simulation.\\~\\
\subsection*{add\_celestial\_bodies(sim\_data: Dict[str, Any]) -> None:}\nAdds celestial body structures (for Sun and Moon) to the simulation data.\\~\\
\subsection*{add\_fixed\_points(sim\_data: Dict[str, Any], num\_points: int = 100) -> None:}\nAdds a structure for fixed reference points in the GCRS frame.

\section*{visibility.py}
\subsection*{solarexclusion(data\_struct: Dict[str, Any])}
Calculates solar exclusion for all satellites based on their pointing vectors.
\subsection*{exclusion(data\_struct: Dict[str, Any], satellite\_index: int, print\_debug: bool = False)}
Determines if a satellite's pointing vector is excluded by the Sun, Moon, or Earth.
\subsection*{update\_visibility\_table(data\_struct: Dict[str, Any], print\_debug\_for\_sat: Optional[int] = None)}
Updates the visibility table for each satellite against each fixed point.

\end{document}
